\section{Architecture}


\subsection{Tensor Core}

\begin{figure}[t]
\centering
\begin{tikzpicture}[
    node distance=0.8cm,
    box/.style={rectangle, draw, minimum width=2.5cm, minimum height=0.7cm, align=center, font=\small},
    arrow/.style={->, thick}
]
\node[box, fill=blue!10] (tensor) {Tensor};
\node[box, below=of tensor] (storage) {Storage};
\node[box, below left=0.5cm and 0cm of storage] (cpu) {CpuStorage};
\node[box, below=of storage] (cuda) {CudaStorage};
\node[box, below right=0.5cm and 0cm of storage] (metal) {MetalStorage};
\node[box, right=2cm of tensor, fill=green!10] (op) {BackpropOp};
\node[box, below=of op] (layout) {Layout};

\draw[arrow] (tensor) -- (storage);
\draw[arrow] (storage) -- (cpu);
\draw[arrow] (storage) -- (cuda);
\draw[arrow] (storage) -- (metal);
\draw[arrow] (tensor) -- (op);
\draw[arrow] (tensor) -- (layout);
\end{tikzpicture}
\caption{Tensor architecture with pluggable storage backends.}
\label{fig:tensor}
\end{figure}

The core \texttt{Tensor} type (Figure~\ref{fig:tensor}) wraps:

\begin{lstlisting}[language=Rust]
pub struct Tensor_(Arc<Tensor_>);

struct Tensor_ {
    id: TensorId,
    storage: Arc<RwLock<Storage>>,
    layout: Layout,
    op: BackpropOp,
    is_variable: bool,
    dtype: DType,
    device: Device,
}
\end{lstlisting}

\paragraph{Reference Counting.} Tensors use \texttt{Arc} for cheap cloning during graph construction. Storage is independently reference-counted to enable view operations without copying data.

\paragraph{Layout.} The \texttt{Layout} struct tracks shape, strides, and offset, enabling non-contiguous views:

\begin{lstlisting}[language=Rust]
pub struct Layout {
    shape: Shape,
    stride: Vec<usize>,
    start_offset: usize,
}
\end{lstlisting}

\paragraph{Backpropagation.} The \texttt{BackpropOp} enum records the operation graph for automatic differentiation:

\begin{lstlisting}[language=Rust]
pub enum Op {
    Binary(Tensor, Tensor, BinaryOp),
    Unary(Tensor, UnaryOp),
    Matmul(Tensor, Tensor),
    Reduce(Tensor, ReduceOp, Vec<usize>),
    Conv2D { arg, kernel, padding, stride, dilation },
    // ... 30+ operations
}
\end{lstlisting}

\subsection{Data Types}

The \texttt{DType} enum supports:

\begin{lstlisting}[language=Rust]
pub enum DType {
    // Standard floats
    F16, BF16, F32, F64,
    // Integers
    U8, U32, I32, I64,
    // Experimental FHE-friendly types
    F8E4M3,  // 8-bit float
    F4,      // 4-bit float
    F6E2M3, F6E3M2,  // 6-bit floats
}
\end{lstlisting}

The \texttt{WithDType} trait provides type-level operations:

\begin{lstlisting}[language=Rust]
pub trait WithDType: Sized + Copy {
    const DTYPE: DType;
    fn to_cpu_storage(data: &[Self]) -> CpuStorage;
    fn from_cpu_storage(s: &CpuStorage) -> Vec<Self>;
}
\end{lstlisting}

\subsection{Backend System}

Three backends implement the \texttt{BackendStorage} trait:

\paragraph{CPU Backend.} Pure Rust implementation with SIMD optimizations:

\begin{lstlisting}[language=Rust]
#[cfg(target_arch = "x86_64")]
mod avx {
    use std::arch::x86_64::*;

    pub fn vec_add_f32(a: &[f32], b: &[f32],
                       c: &mut [f32]) {
        for i in (0..a.len()).step_by(8) {
            unsafe {
                let va = _mm256_loadu_ps(&a[i]);
                let vb = _mm256_loadu_ps(&b[i]);
                let vc = _mm256_add_ps(va, vb);
                _mm256_storeu_ps(&mut c[i], vc);
            }
        }
    }
}
\end{lstlisting}

Runtime dispatch selects optimal kernels based on CPU features.

\paragraph{CUDA Backend.} Uses cuBLAS for matrix operations and custom kernels for element-wise ops. Optional cuDNN integration for convolutions.

\paragraph{Metal Backend.} Apple Silicon support via Metal Performance Shaders and custom compute kernels for M1/M2/M3 optimization.

\subsection{Operations}

\subsubsection{Unary Operations}

Defined via macro for consistency:

\begin{lstlisting}[language=Rust]
macro_rules! unary_op {
    ($fn_name:ident, $op_name:ident) => {
        pub fn $fn_name(&self) -> Result<Self> {
            let storage = self.storage()
                .unary_impl::<op::$op_name>(
                    self.layout())?;
            let op = BackpropOp::new1(self,
                |s| Op::Unary(s, UnaryOp::$op_name));
            Ok(from_storage(storage, shape, op, false))
        }
    };
}

unary_op!(exp, Exp);
unary_op!(log, Log);
unary_op!(sin, Sin);
unary_op!(gelu, Gelu);
unary_op!(silu, Silu);
\end{lstlisting}

\subsubsection{Binary Operations}

Support broadcasting with shape validation:

\begin{lstlisting}[language=Rust]
pub fn broadcast_add(&self, rhs: &Self) -> Result<Self> {
    let shape = self.shape()
        .broadcast_shape_binary_op(rhs.shape(), "add")?;
    // ... broadcast and compute
}
\end{lstlisting}

\subsubsection{Reductions}

Sum, mean, max, min, argmax, argmin with axis specification:

\begin{lstlisting}[language=Rust]
let x = Tensor::randn((2, 3, 4), DType::F32, &device)?;
let sum_all = x.sum_all()?;           // Scalar
let sum_axis = x.sum(1)?;             // Shape: (2, 4)
let max_keepdim = x.max_keepdim(2)?;  // Shape: (2, 3, 1)
\end{lstlisting}

